%───────────────────────────────────────────────────────────────────────────────
\chapter{Matemáticas semánticas}
\label{ch:ecuaciones}
%───────────────────────────────────────────────────────────────────────────────

La tipografía matemática estándar codifica una sola dimensión de información:
el tipo matemático. Los vectores van en negrita, los operadores en redonda,
los escalares en itálica. Pero una ecuación contiene al menos tres dimensiones
de información simultáneas que la tipografía convencional no distingue: el
\textit{rol sintáctico} del símbolo (¿es un operador, una función, un dominio?),
su \textit{estatus epistémico} (¿es una incógnita, un dato fijo, una constante
universal?) y su \textit{tipo matemático} (¿es un vector, un tensor, un escalar?).

El sistema de coloreado semántico de este template codifica estas tres
dimensiones mediante color, con un mapa fijo entre rol y tono que es
consistente en todo el documento.

%───────────────────────────────────────────────────────────────────────────────
\section{Regla fundamental: siempre en modo matemático}
\label{sec:regla-matmode}
%───────────────────────────────────────────────────────────────────────────────

\begin{definicion}{Modo matemático obligatorio}
\label{def:matmode}
Todas las macros de coloreado semántico (\texttt{\textbackslash eqop},
\texttt{\textbackslash equ}, \texttt{\textbackslash eqk}, etc.) deben usarse
\textbf{dentro de un entorno matemático}. Fuera de modo matemático, los
símbolos que colorean son glifos de texto, no glifos matemáticos: el color
se aplica sobre el carácter incorrecto o no se aplica en absoluto.
\end{definicion}

La forma correcta e incorrecta es:

\begin{lstlisting}[style=latex, frame=none, numbers=none]
% CORRECTO: dentro de modo matematico
$\eqop{\nabla} \equ{f}$          % inline
\[ \eqop{\nabla}^2 \equ{\phi} = \eqk{\rho} / \eqcon{\varepsilon_0} \]

% INCORRECTO: fuera de modo matematico
\eqop{\nabla} \equ{f}   % el color se aplica sobre texto, no glifo matematico
\end{lstlisting}

Esto aplica en todos los contextos: texto corrido, tablas, títulos de
entornos, labels de ejes en pgfplots, labels de nodos en TikZ. En cada
caso, el wrapper matemático es obligatorio.

%───────────────────────────────────────────────────────────────────────────────
\section{Los tres ejes de coloreado}
\label{sec:tres-ejes}
%───────────────────────────────────────────────────────────────────────────────

\dificultad{2}{3}

\subsection{Eje 1 — Rol sintáctico}
\label{subsec:eje-sintactico}

El rol sintáctico de un símbolo es estable entre contextos: $\nabla$ es
siempre un operador diferencial, independientemente de qué ecuación aparezca
o cuál sea la incógnita.

\begin{acadtable}[H]{tab:eje-sintactico}{Eje 1: Rol sintáctico}
  \begin{tabularx}{\textwidth}{l l l X}
    \toprule
    \headerrow
    \thead{Macro} & \thead{Alias} & \thead{Color} & \thead{Cubre} \\
    \midrule
    \texttt{\textbackslash eqop\{·\}}
      & \texttt{eq-op} & \textcolor{eq-op}{\textbf{Blue}}
      & Operadores: $\eqop{\int}$, $\eqop{\sum}$, $\eqop{\partial}$,
        $\eqop{\nabla}$, $\eqop{\times}$, $\eqop{\delta}(\cdot,\cdot)$ \\
    \texttt{\textbackslash eqfn\{·\}}
      & \texttt{eq-fn} & \textcolor{eq-fn}{\textbf{Green}}
      & Funciones/mapas: $\eqfn{f}(\cdot)$, $\eqfn{\sigma}(\cdot)$,
        $\eqfn{V}(\mathbf{r})$, $\eqfn{\mathcal{L}}$ \\
    \texttt{\textbackslash eqdm\{·\}}
      & \texttt{eq-dm} & \textcolor{eq-dm}{\textbf{Teal}}
      & Diferenciales/medidas: $\eqdm{dx}$, $\eqdm{dt}$,
        $\eqdm{dW}$, $\eqdm{d^3r}$, $\eqdm{d\mu}$ \\
    \texttt{\textbackslash eqdom\{·\}}
      & \texttt{eq-dom} & \textcolor{eq-dom}{\textbf{Sapphire}}
      & Dominios/conjuntos: $\eqdom{\Omega}$, $\eqdom{\partial\Omega}$,
        $\eqdom{\mathbb{R}^n}$, $\eqdom{\Sigma}$, $\eqdom{Q}$ \\
    \bottomrule
  \end{tabularx}
\end{acadtable}

\subsection{Eje 2 — Estatus epistémico}
\label{subsec:eje-epistemico}

El estatus epistémico es la dimensión que la tipografía convencional
\textit{no} codifica. En $f(\mathbf{x};\boldsymbol{\theta})$, $\mathbf{x}$
es dato de entrada y $\boldsymbol{\theta}$ es el parámetro a aprender ---
en inferencia estadística es exactamente al revés. El color hace explícita
esta inversión sin texto adicional.

\begin{acadtable}[H]{tab:eje-epistemico}{Eje 2: Estatus epistémico}
  \begin{tabularx}{\textwidth}{l l l X}
    \toprule
    \headerrow
    \thead{Macro} & \thead{Alias} & \thead{Color} & \thead{Cubre} \\
    \midrule
    \texttt{\textbackslash equ\{·\}}
      & \texttt{eq-unknown} & \textcolor{eq-unknown}{\textbf{Peach}}
      & Incógnita / buscado / aprendido: la cantidad que se desea hallar \\
    \texttt{\textbackslash eqk\{·\}}
      & \texttt{eq-known} & \textcolor{eq-known}{\textbf{Mauve}}
      & Conocido / dato / parámetro fijado: entra al sistema como input \\
    \texttt{\textbackslash eqcon\{·\}}
      & \texttt{eq-const} & \textcolor{eq-const}{\textbf{Lavender}}
      & Constante universal/estructural: $\eqcon{\hbar}$, $\eqcon{k_B}$,
        $\eqcon{\pi}$, $\eqcon{\varepsilon}$ \\
    \bottomrule
  \end{tabularx}
\end{acadtable}

\subsection{Eje 3 — Tipo matemático}
\label{subsec:eje-tipo}

El eje de tipo refuerza lo que la tipografía ya comunica: la negrita indica
vector, la doble barra indica espacio. El color añade un canal redundante
útil en documentos donde la distinción rango-tensor es semánticamente
importante y el lector necesita identificarla de un vistazo.

\begin{acadtable}[H]{tab:eje-tipo}{Eje 3: Tipo matemático}
  \begin{tabularx}{\textwidth}{l l l X}
    \toprule
    \headerrow
    \thead{Macro} & \thead{Alias} & \thead{Color} & \thead{Cubre} \\
    \midrule
    \texttt{\textbackslash eqvec\{·\}}
      & \texttt{eq-vec} & \textcolor{eq-vec}{\textbf{Red}}
      & Vectores/campos: $\eqvec{\mathbf{u}}$, $\eqvec{\mathbf{E}}$,
        $\eqvec{\mathbf{B}}$, $\eqvec{\mathbf{J}}$ \\
    \texttt{\textbackslash eqten\{·\}}
      & \texttt{eq-ten} & \textcolor{eq-ten}{\textbf{Rosewater}}
      & Tensores/matrices: $\eqten{\sigma_{ij}}$, $\eqten{g_{\mu\nu}}$,
        $\eqten{W}$, $\eqten{\Sigma}$ \\
    \bottomrule
  \end{tabularx}
\end{acadtable}

%───────────────────────────────────────────────────────────────────────────────
\section{Ecuaciones completas con los tres ejes}
\label{sec:ecuaciones-completas}
%───────────────────────────────────────────────────────────────────────────────

La riqueza del sistema aparece cuando los tres ejes se aplican
simultáneamente en una ecuación completa. Los ejemplos siguientes
usan física y machine learning para mostrar el sistema bajo carga real.

\subsection{Mecánica cuántica: ecuación de Schrödinger}
\label{subsec:schrodinger}

\begin{fisicaconexion}{Ecuación de Schrödinger~\cite{griffiths2017quantum}}
La ecuación de Schrödinger dependiente del tiempo:
\[
  \eqop{i}\eqcon{\hbar}\,\eqop{\frac{\partial}{\partial t}}\equ{\Psi}
  = \eqop{\hat{H}}\equ{\Psi}
  = \left(\eqop{-}\frac{\eqcon{\hbar}^2}{2\eqk{m}}\eqop{\nabla}^2
    + \eqk{V}(\eqvec{\mathbf{r}},t)\right)\equ{\Psi}
\]
El color comunica: $\eqcon{\hbar}$ y $\eqcon{i}$ son constantes estructurales.
$\eqk{m}$ es la masa, un parámetro fijado del sistema.
$\eqk{V}(\mathbf{r},t)$ es el potencial, conocido.
$\equ{\Psi}(\mathbf{r},t)$ es la función de onda, la incógnita que se busca.
$\eqop{\partial/\partial t}$ y $\eqop{\nabla^2}$ son operadores diferenciales.
\end{fisicaconexion}

\subsection{Electrodinámica: ecuaciones de Maxwell}
\label{subsec:maxwell}

\begin{fisicaconexion}{Ecuaciones de Maxwell~\cite{jackson1999electrodynamics}}
En forma diferencial, con unidades SI:
\begin{align}
  \eqop{\nabla \cdot} \eqvec{\mathbf{E}}
    &= \frac{\eqk{\rho}}{\eqcon{\varepsilon_0}}
  \label{eq:gauss-e} \\
  \eqop{\nabla \cdot} \eqvec{\mathbf{B}}
    &= 0
  \label{eq:gauss-b} \\
  \eqop{\nabla \times} \eqvec{\mathbf{E}}
    &= \eqop{-}\eqop{\frac{\partial}{\partial t}}\eqvec{\mathbf{B}}
  \label{eq:faraday} \\
  \eqop{\nabla \times} \eqvec{\mathbf{B}}
    &= \eqcon{\mu_0}\left(\eqvec{\mathbf{J}}
      + \eqcon{\varepsilon_0}\eqop{\frac{\partial}{\partial t}}
      \eqvec{\mathbf{E}}\right)
  \label{eq:ampere}
\end{align}
Los campos $\eqvec{\mathbf{E}}$ y $\eqvec{\mathbf{B}}$ son incógnitas
vectoriales en rojo. Las fuentes $\eqk{\rho}$ y $\eqvec{\mathbf{J}}$ son
datos del problema en malva. Las constantes $\eqcon{\varepsilon_0}$ y
$\eqcon{\mu_0}$ en lavanda son universales.
\end{fisicaconexion}

\subsection{Machine learning: ELBO variacional}
\label{subsec:elbo}

\begin{ejemplo}{Evidence Lower BOund (ELBO)~\cite{goodfellow2016deep}}
\label{ex:elbo}
En inferencia variacional, el objetivo es maximizar:
\[
  \eqfn{\mathcal{L}}(\eqk{\phi}, \eqk{\theta})
  = \eqop{\mathbb{E}}_{\equ{z} \sim \eqfn{q}_{\eqk{\phi}}}
    \!\left[\log \eqfn{p}_{\eqk{\theta}}(\eqk{x} \mid \equ{z})\right]
  \eqop{-} \eqop{D_{\mathrm{KL}}}\!\left(
    \eqfn{q}_{\eqk{\phi}}(\equ{z} \mid \eqk{x})
    \,\|\, \eqfn{p}(\equ{z})
  \right)
\]
El color hace explícita la estructura: $\eqk{\phi}$ y $\eqk{\theta}$ son los
parámetros del encoder y decoder, datos del sistema de optimización.
$\equ{z}$ es la variable latente, la incógnita que se integra.
$\eqfn{q}_\phi$ y $\eqfn{p}_\theta$ son las distribuciones (funciones del
modelo). $\eqop{\mathbb{E}}$ y $\eqop{D_{\mathrm{KL}}}$ son operadores.
\end{ejemplo}

\subsection{Física computacional: ecuación de Fokker-Planck}
\label{subsec:fokker-planck}

\begin{definicion}{Ecuación de Fokker-Planck}
\label{def:fokker-planck}
La evolución de la densidad de probabilidad $\equ{p}(\eqvec{\mathbf{x}},t)$
de un proceso de difusión gobernado por la SDE
$\eqdm{d}\equ{X} = \eqfn{f}(\equ{X})\eqdm{dt} + \eqk{\sigma}\eqdm{dW}$
está dada por:
\[
  \eqop{\frac{\partial}{\partial t}}\equ{p}
  = \eqop{-}\eqop{\nabla_x \cdot}
    \!\left(\eqfn{f}\,\equ{p}\right)
  + \frac{\eqk{\sigma}^2}{2}\,\eqop{\nabla_x^2}\,\equ{p}
\]
donde $\eqop{\nabla_x \cdot}$ es el operador divergencia y
$\eqop{\nabla_x^2}$ es el Laplaciano, ambos actuando sobre la
variable espacial $\eqvec{\mathbf{x}}$. El coeficiente de difusión
$\eqk{\sigma}^2/2$ es un parámetro conocido.
\end{definicion}

%───────────────────────────────────────────────────────────────────────────────
\section{Macros compuestos}
\label{sec:macros-compuestos}
%───────────────────────────────────────────────────────────────────────────────

\texttt{config/equations.tex} define macros compuestos para patrones
frecuentes en física y machine learning. Usan los ejes semánticos
internamente.

\begin{acadtable}[H]{tab:macros-compuestos}{Macros compuestos de equations.tex}
  \begin{tabularx}{\textwidth}{l X X}
    \toprule
    \headerrow
    \thead{Macro} & \thead{Uso} & \thead{Produce} \\
    \midrule
    \texttt{\textbackslash eqpdv\{f\}\{t\}}
      & Derivada parcial coloreada
      & $\eqpdv{f}{t}$ \\
    \texttt{\textbackslash eqgrad\{u\}}
      & Gradiente coloreado
      & $\eqgrad{u}$ \\
    \texttt{\textbackslash eqdvg\{\textbackslash mathbf\{u\}\}}
      & Divergencia coloreada
      & $\eqdvg{\mathbf{u}}$ \\
    \texttt{\textbackslash eqSDE\{X\}\{f(X)\}\{\textbackslash sigma\}}
      & SDE coloreada completa
      & $\eqSDE{X}{f(X)}{\sigma}$ \\
    \bottomrule
  \end{tabularx}
\end{acadtable}

El macro \texttt{\textbackslash eqpdv} aplica la convención epistémica
natural de las derivadas parciales: el numerador es siempre la incógnita
(lo que cambia, lo que se busca entender) y el denominador es la variable
independiente, dato del problema:

\begin{lstlisting}[style=latex, frame=none, numbers=none]
% Ecuacion de calor: u es la incognita, t y x son variables independientes
\[
  \eqpdv{u}{t} = \eqk{\alpha}\, \eqpdv[2]{u}{x}
\]
\end{lstlisting}

\[
  \eqpdv{u}{t} = \eqk{\alpha}\, \eqop{\frac{\partial^2 \equ{u}}{\partial \eqk{x}^2}}
\]

%───────────────────────────────────────────────────────────────────────────────
\section{Coloreado en contextos distintos}
\label{sec:contextos}
%───────────────────────────────────────────────────────────────────────────────

El sistema funciona en cualquier contexto donde se pueda escribir matemáticas.
La regla del modo matemático (\cref{def:matmode}) aplica en todos.

\subsection{En tablas}
\label{subsec:eq-en-tablas}

El documento de referencia muestra el caso de expresiones regulares con
theming en tablas: $\eqk{a} \eqop{\mid} \eqk{b}$ en una celda de
\texttt{tabularx}. El patrón es simplemente matemáticas normales dentro de la celda:

\begin{acadtable}[H]{tab:regexp-coloreado}{Expresiones regulares con coloreado semántico}
  \begin{tabularx}{\textwidth}{X X}
    \toprule
    \headerrow
    \thead{Expresión regular} & \thead{Lenguaje descrito} \\
    \midrule
    $\eqk{a} \eqop{\mid} \eqk{b}$
      & $\eqdom{\{}\eqk{a}, \eqk{b}\eqdom{\}}$ \\
    $\eqk{a}\eqop{*}$
      & $\eqdom{\{}\eqcon{\varepsilon}, \eqk{a}, \eqk{aa}, \ldots\eqdom{\}}$ \\
    $(\eqk{a} \eqop{\mid} \eqk{b})\eqop{*}\eqk{a}\eqk{b}\eqk{b}$
      & Cadenas sobre $\eqdom{\{a,b\}}$ con sufijo $\eqk{abb}$ \\
    $\eqop{\delta} : \eqdom{Q} \eqop{\times} \eqdom{\Sigma}
      \eqop{\rightarrow} \eqdom{Q}$
      & Función de transición de un DFA \\
    \bottomrule
  \end{tabularx}
\end{acadtable}

\subsection{En algoritmos}
\label{subsec:eq-en-algoritmos}

Dentro de \texttt{algorithm2e}, el theming funciona directamente en modo
matemático. El documento de referencia muestra algoritmos de construcción
de subconjuntos con $\eqk{q'_0} \eqop{\gets} \eqop{E}(\eqdom{\{}\eqk{q_0}\eqdom{\}})$ ---
el estado inicial (conocido) se inicializa con el operador de clausura:

\begin{lstlisting}[style=latex, frame=none, numbers=none]
\begin{algorithm}[H]
  \KwIn{NFA $\eqk{N} = (\eqdom{Q}, \eqdom{\Sigma}, \eqop{\delta},
               \eqk{q_0}, \eqdom{F})$}
  $\eqk{q'_0} \eqop{\gets} \eqop{E}(\eqdom{\{}\eqk{q_0}\eqdom{\}})$\;
  \While{$\equ{T} \eqop{\in} \eqdom{Q'}$ no marcado}{
    $\equ{U} \eqop{\gets} \eqop{E}(\eqop{move}(\equ{T}, \equ{a}))$\;
  }
\end{algorithm}
\end{lstlisting}

\subsection{En TikZ}
\label{subsec:eq-en-tikz}

En nodos y aristas de TikZ, los labels matemáticos van dentro de \texttt{\$...\$}.
El theming es idéntico al de cualquier ecuación:

\begin{lstlisting}[style=latex, frame=none, numbers=none]
\path[->]
  (q0) edge node {$\eqk{a}$} (q1)
  (q0) edge node {$\eqcon{\varepsilon}$} (q2);
\end{lstlisting}

\begin{observacion}{Colores en TikZ: usar aliases, no Ctp*}
El documento de referencia usa \texttt{CtpBlue}, \texttt{CtpBase},
\texttt{CtpText} directamente en los estilos de nodos TikZ.
Esto viola la filosofía del template: los colores \texttt{Ctp*}
deben vivir \textit{únicamente} en \texttt{themes/catppuccin-palette.tex}.
El uso correcto es:
\begin{lstlisting}[style=latex, frame=none, numbers=none]
% INCORRECTO: Ctp* en un archivo de contenido
state/.style={draw=CtpBlue, fill=CtpBase, text=CtpText}

% CORRECTO: aliases semanticos
state/.style={draw=tikz-node-border, fill=background, text=maintext}
\end{lstlisting}
\end{observacion}

%───────────────────────────────────────────────────────────────────────────────
\section{Toggle de impresión en blanco y negro}
\label{sec:toggle-bn}
%───────────────────────────────────────────────────────────────────────────────

\begin{definicion}{Toggle \texttt{eqcolors}}
\label{def:toggle}
El toggle global \texttt{\textbackslash ifeqcolors} controla si las macros
aplican color o pasan el argumento sin modificar. Esto permite generar una
versión B/N del documento sin editar ninguna ecuación:
\end{definicion}

\begin{lstlisting}[style=latex, frame=none, numbers=none]
% En config/equations.tex (activo por defecto):
\eqcolorstrue    % color activado
\eqcolorsfalse   % color desactivado: todas las macros son transparentes

% Para desactivar en un grupo local:
{\eqcolorsfalse
  \begin{equation}
    i\hbar \frac{\partial \Psi}{\partial t} = \hat{H}\Psi
  \end{equation}
}
\end{lstlisting}

Con \texttt{\textbackslash eqcolorsfalse}, \texttt{\textbackslash eqop\{\textbackslash nabla\}}
produce $\nabla$ sin color --- funcionalmente idéntico a escribir
$\nabla$ directamente. El documento compila sin cambios.

%───────────────────────────────────────────────────────────────────────────────
\section{Integración con el paquete \texttt{physics}}
\label{sec:physics-pkg}
%───────────────────────────────────────────────────────────────────────────────

El paquete \texttt{physics}~\cite{physics-pkg} provee comandos de alto nivel
que se integran naturalmente con el coloreado semántico. La norma, el valor
absoluto y los operadores de Dirac se colorean envolviendo el comando de
\texttt{physics} en la macro semántica correcta:

\begin{lstlisting}[style=latex, frame=none, numbers=none]
\norm{\equ{f}}_{\eqdom{L^2}}         % norma de la incognita en L2
\abs{\eqk{z} - \eqk{z_0}}            % valor absoluto de datos
\bra{\eqk{\phi}}\hat{H}\ket{\equ{\psi}}   % elemento de matriz
\end{lstlisting}

\[
  \|\equ{f}\|_{\eqdom{L^2}}^2
  = \eqop{\int}_{\eqdom{\Omega}} \abs{\equ{f}(\eqvec{\mathbf{x}})}^2
    \eqdm{d\mu}(\eqvec{\mathbf{x}})
\]

\begin{observacion}{Conflicto de \texttt{\textbackslash div}}
El paquete \texttt{physics} redefine \texttt{\textbackslash div} como
el operador divergencia ($\nabla\cdot$). Para la divergencia coloreada
usar \texttt{\textbackslash eqdvg\{\}}, definido en
\texttt{config/equations.tex} precisamente para evitar este conflicto.
Para la división aritmética usar \texttt{\textbackslash frac\{\}\{\}} o
el operador \texttt{/}.
\end{observacion}

%───────────────────────────────────────────────────────────────────────────────
\printbibliography[heading=subbibliography]
